\hypertarget{1-question}{%
\subsection{1. Question}\label{1-question}}

Can an accept-or-abstain gate, calibrated on data held out \emph{by
chemistry}, keep the precision of a molecular model\textquotesingle s
accepted predictions at a stated target when the queries come from
chemistry the model has not seen? And does combining an input-distance
signal with a model-confidence signal rank errors better than either
signal alone?

\hypertarget{2-definitions-fixed-in-advance}{%
\subsection{2. Definitions fixed in
advance}\label{2-definitions-fixed-in-advance}}

\begin{itemize}
\tightlist
\item
  \textbf{Correct prediction.} LogP: \textbar y − ŷ\textbar{} ≤ 0.5 log
  units. Reaction yield: \textbar y − ŷ\textbar{} ≤ 10 percentage
  points.
\item
  \textbf{Coverage} is the share of queries the gate accepts.
  \textbf{Selective precision} is the share of accepted queries that are
  correct. \textbf{Selective error rate} = 1 − selective precision.
\item
  \textbf{Target precisions:} t = 0.8 and t = 0.9. Risk level for the
  threshold guarantee: δ = 0.1.
\end{itemize}

\hypertarget{3-data-all-public}{%
\subsection{3. Data (all public)}\label{3-data-all-public}}

\begin{itemize}
\tightlist
\item
  \textbf{D1 LogP pool:} SangsterLogP main sheet (Zenodo
  10.5281/zenodo.19387551, CC BY 4.0). Parsed with RDKit; unparsable
  SMILES dropped; duplicates by InChIKey merged (first record kept).
  Every compound whose InChIKey occurs in either external set below is
  removed from the pool.
\item
  \textbf{E1, prospective external set:} the SangsterLogP "Prospective"
  sheet (327 compounds).
\item
  \textbf{E2, fluorinated external set:} the 183 measured amide logP
  values of \texttt{blackthorn-ai/fluoricLogPpKa} (MIT), each
  fluorinated compound beside its non-fluorinated parent.
\item
  \textbf{D3a Buchwald--Hartwig} (Ahneman et al. 2018; as redistributed
  in \texttt{SoftServeInc/yield-paper}, MIT): 3,955 reactions; the four
  published out-of-sample additive tests (sheets Test1--Test4, training
  rows 1--3057, 1--3055, 1--3058 and 1--3055 respectively) and random
  splits.
\item
  \textbf{D3b Suzuki--Miyaura} (Perera et al. 2018; same
  redistribution): 5,760 reactions, target
  \texttt{Product\_Yield\_PCT\_Area\_UV}.
\end{itemize}

\hypertarget{4-representations-and-base-models}{%
\subsection{4. Representations and base
models}\label{4-representations-and-base-models}}

\begin{itemize}
\tightlist
\item
  Molecules: ECFP4 counts (radius 2, 2,048 bins) plus ten RDKit
  descriptors (MolWt, TPSA, NumHDonors, NumHAcceptors,
  NumRotatableBonds, RingCount, NumAromaticRings, FractionCSP3,
  HeavyAtomCount, NumHeteroatoms). Crippen MolLogP is deliberately
  excluded.
\item
  Reactions: concatenated ECFP4 bit vectors (1,024 bits) of each
  variable component (BH: aryl halide, ligand, base, additive; SM:
  reactant 1, reactant 2, ligand, reagent, solvent). Component
  structures follow the cited redistribution.
\item
  \textbf{M1 (primary):} random forest, 300 trees, max\_features = 0.2,
  other settings at scikit-learn defaults. Spread u = standard deviation
  of the per-tree predictions.
\item
  \textbf{M2 (secondary, robustness only):} ensemble of five multilayer
  perceptrons (512--128, ReLU, Adam, L2 = 1e-4, early stopping on 10\%
  of the fit set). Spread u = standard deviation across members.
\item
  Hyper-parameters are fixed here and are not tuned on any calibration
  or test fold.
\end{itemize}

\hypertarget{5-splits-five-seeds-04}{%
\subsection{5. Splits (five seeds, 0--4)}\label{5-splits-five-seeds-04}}

For each dataset a pool is divided by a group label: Butina clusters of
ECFP4 at Tanimoto distance 0.65 (D1), additive (D3a), ligand (D3b).

\begin{itemize}
\tightlist
\item
  \textbf{Held-out test H:} D1: 15\% of clusters; D3a: the published
  TestN additives; D3b: two ligands per seed.
\item
  \textbf{Held-out calibration C\_h:} groups disjoint from H and from
  the fit set. D1: 15\% of clusters; D3a: four additives drawn from the
  TestN training additives; D3b: two further ligands.
\item
  \textbf{In-distribution part I} (all remaining groups) is split at
  random into the fit set F (70\%), the random calibration fold C\_r
  (15\%) and the random test R (15\%).
\item
  The base model is trained on F only. Both gates (C\_r-calibrated and
  C\_h-calibrated) use the same model.
\item
  \textbf{D1 external tests:} the D1 models are also scored on E1. A
  separate D1 model, trained after removing every compound with an sp³
  C--F bond (SMARTS \texttt{{[}CX4{]}F}) from the pool, is scored on E2;
  fluorinated compounds and parents are reported apart.
\end{itemize}

Held-out tests are therefore: D1 cluster, D1--E1, D1--E2 (D1 family);
D3a Test1--Test4 and D3b ligand (D3 family).

\hypertarget{6-signals}{%
\subsection{6. Signals}\label{6-signals}}

\begin{itemize}
\tightlist
\item
  \textbf{s1:} maximum Tanimoto similarity of the query fingerprint to
  F; \textbf{s5:} mean of the five largest.
\item
  \textbf{u:} ensemble spread of the base model.
\item
  \textbf{Conformal intervals (nominal 90\%):} normalised split
  conformal, score \textbar y − ŷ\textbar{} / (u + ε), ε = 0.1 ×
  median(u) on the calibration fold.
\end{itemize}

\hypertarget{7-gates}{%
\subsection{7. Gates}\label{7-gates}}

\begin{itemize}
\tightlist
\item
  Each calibration fold C is split at random into halves C\_a and C\_b.
\item
  \textbf{Single-signal gates:} score s1, or −u.
\item
  \textbf{Gate B (fused):} logistic regression of correctness on (s1,
  s5, log u), fitted on C\_a. Its score is the predicted probability of
  a correct prediction.
\item
  \textbf{Threshold by Learn-then-Test} on C\_b, for every score:
  candidate thresholds are the score quantiles of C\_b, tested in fixed
  sequence from the strictest that accepts at least 50 points of C\_b,
  loosening one step at a time. A threshold passes when the binomial
  p-value P(Bin(n, t) ≥ k) ≤ δ, where n points are accepted and k are
  correct. The loosest threshold reached before the first failure is
  used. If none passes, the gate abstains on every query.
\item
  \textbf{Gate A (conjunctive, no guarantee):} accept if s1 ≥ $\tau_1$ and u ≤
  $\tau_2$; ($\tau_1$, $\tau_2$) from a 20 × 20 grid of quantiles of the whole C, chosen
  as the largest coverage whose empirical precision on C is at least t.
\end{itemize}

\hypertarget{8-metrics}{%
\subsection{8. Metrics}\label{8-metrics}}

Per test set, model and seed: RMSE, MAE and R² ungated; error rate
ungated; risk--coverage curves of selective error rate and selective
RMSE for each score; AURC of the selective error rate (mean over all
coverage levels k/n); coverage and achieved precision of every gate at t
= 0.8 and 0.9; empirical coverage of the conformal intervals on all
queries and on accepted queries. Results are averaged over seeds; 95\%
intervals come from 1,000 bootstrap resamples of test queries, pooled
over seeds.

\hypertarget{9-hypotheses-and-pass-criteria-model-m1}{%
\subsection{9. Hypotheses and pass criteria (model
M1)}\label{9-hypotheses-and-pass-criteria-model-m1}}

\begin{itemize}
\tightlist
\item
  \textbf{H1, the gate helps where the model is weak.} On every held-out
  test, the selective error rate of the fused score at 30\% coverage is
  lower than the ungated error rate, with the 95\% bootstrap interval of
  the difference excluding zero.
\item
  \textbf{H2, the calibration split matters.} Over the (held-out test ×
  target) cells, Gate B calibrated on C\_r falls short of the target by
  more than 0.05 in at least half of the cells; calibrated on C\_h, it
  reaches at least t − 0.05 in at least three quarters of the cells
  (seed-averaged precision).
\item
  \textbf{H3, the combination adds something.} The AURC of Gate
  B\textquotesingle s score (C\_h-fitted) is at least 10\% lower,
  relatively, than the better single signal on at least two held-out
  tests in each of the D1 and D3 families.
\item
  \textbf{H4, conformal coverage under shift.} Intervals calibrated on
  C\_r cover at least 0.88 of R; cover less than 0.85 on at least one
  held-out test; and, on each held-out test where they fall below 0.85,
  cover at least 0.88 of the queries accepted by Gate B (C\_h, t = 0.8).
\end{itemize}

Every hypothesis is reported whether or not it holds.

\hypertarget{10-exclusions}{%
\subsection{10. Exclusions}\label{10-exclusions}}

No client, proprietary or confidential data, code or results are used.
Biological context (gene expression, perturbation or cell-line shift)
and generative models are out of scope.

\hypertarget{notes-on-implementation}{%
\subsection{Notes on implementation}\label{notes-on-implementation}}

\begin{itemize}
\tightlist
\item
  A gate that abstains on every query has no defined precision. Such
  seeds are left out of the seed average. A cell in which every seed
  abstains counts as meeting the target in H2.
\item
  H4 cannot be checked on a test where Gate B accepts nothing; this
  counts as a failure of H4 for that test.
\item
  The held-out tests for H1--H4 are the logP cluster test, E1, E2 (all
  183 compounds), BH Test1--Test4 and the SM ligand test.
\item
  Holding out four calibration additives leaves about 1,640 reactions in
  F for each BH test. Ungated BH accuracy is therefore not directly
  comparable with published results on the same tests.
\item
  The logP held-out sets were formed by assigning whole clusters until
  15\% of the compounds was reached.
\item
  The Suzuki--Miyaura component structures come from a table written for
  this study (\texttt{src/aoa/data.py}). The table omits spectator
  fragments: the internal standard, HCl and water of hydration.
\item
  For test sets that are identical in every seed (E1, E2, BH
  Test1--Test4), bootstrap intervals use the same resample for all
  seeds.
\end{itemize}

\hypertarget{amendment-20-12-november-2026-a-graph-network-as-a-third-learner-m3}{%
\subsection{Amendment 2.0, 12 November 2026: a graph network as a third
learner
(M3)}\label{amendment-20-12-november-2026-a-graph-network-as-a-third-learner-m3}}

\textbf{Fixed before any M3 prediction was computed.} Sections 1--10 and
notes 1.1--1.4 are not changed, and the verdicts of version 1, which M1
decides, stand.

\textbf{Why.} Under version 1 the outcome of H3 differed by learner.
With M1, the fused score ranked errors no better than ensemble spread on
any held-out test (relative AURC reduction between −0.9\% and +1.7\%).
With M2, similarity s1 ranked logP errors better than spread on all
three logP tests, and the fused score lowered the AURC there by
2.8--7.4\%, short of the 10\% criterion; on the reaction tests spread
was the better signal and fusion added nothing. Two fingerprint learners
cannot show whether H3 depends on the model family. A directed
message-passing neural network (D-MPNN) learns its own molecular
representation and is the family most uncertainty benchmarks for
molecular property prediction use. This difference was seen after
version 1 was scored, so the amendment is motivated post hoc; every M3
setting and criterion below is fixed before M3 exists.

\textbf{M3, a graph network.}

\begin{itemize}
\tightlist
\item
  \emph{Software.} Chemprop 2.3.1 (PyTorch 2.15.0, Lightning 2.6.6,
  RDKit 2026.03.6), in a separate environment because it requires an
  older RDKit (versions in \texttt{env/graph-requirements.txt}). Data,
  splits, s1 and s5, gates, intervals and metrics are computed in the
  main environment exactly as for M1 and M2
  (\texttt{scripts/graph\_jobs.py}; the M3 fit is
  \texttt{scripts/graph\_train.py}).
\item
  \emph{Input.} D1: the SMILES of the compound. D3a and D3b: the SMILES
  of the variable components of section 4, joined by "." into one
  multi-fragment graph; an empty component (no ligand or no reagent in
  D3b) is left out.
\item
  \emph{Architecture and training.} Chemprop\textquotesingle s defaults,
  not tuned on any fold: bond-based message passing (hidden size 300,
  depth 3, no dropout), norm aggregation, a feed-forward head with one
  hidden layer of 300, standardised targets, mean-squared-error loss,
  batch size 64 and 50 epochs, with the learning rate rising from 1e-4
  to 1e-3 over two warm-up epochs and decaying to 1e-4. A random 10\% of
  F selects the checkpoint with the lowest validation loss; the other
  90\% trains. Training uses the CPU only, three threads per job.
\item
  \emph{Ensemble.} Five members. Member m of seed s uses the random seed
  100·s + m for its initialisation, data order and validation subset. ŷ
  is the member mean; u is the standard deviation across members
  (population, as for M2).
\item
  \emph{Everything else} is as in sections 3--8: M3 is trained on F
  only, with the same five seeds and splits, and its predictions pass
  through the same gates and metrics.
\end{itemize}

\textbf{Hypotheses for M3.} H1--H4 of section 9 are applied to M3 with
unchanged criteria and read with notes 1.1, 1.3 and 1.4. The question
that motivated this amendment is H3; no direction is predicted.
M3\textquotesingle s verdicts are reported beside M1\textquotesingle s
and M2\textquotesingle s.

\textbf{Descriptive, not tested.} For each held-out test: which single
signal, s1 or −u, has the lower AURC under M3, and the fused
score\textquotesingle s relative AURC change, beside M1 and M2; and the
ungated accuracy of M3 beside M1 and M2. Reporting follows notes 1.3 and
1.4 (per-seed counts; pooled precision is descriptive, with no
interval).

\textbf{Failures.} A member whose predictions are not finite is
retrained once with the seed 100·s + m + 50. If that also fails, the
ensemble keeps its other members and the tables say so. Every task and
seed is reported.

\textbf{Order.} This amendment and the code that produces M3 predictions
are committed before any M3 prediction is computed. Before this
amendment was written, a timing check trained D-MPNN members for two
epochs on 1,000 fit-set rows from three job files (D1, D3a Test1 and
D3b, seed 0), and a smoke test trained a five-member ensemble on 300
fit-set rows of the D3a Test1 job (seed 0) and predicted 50 further
fit-set rows, twice, with identical results. Neither computed a
prediction on any calibration, test or external set.

\textbf{2.1, 12 November 2026, after M3 was scored: documentation only.}
Nothing in amendment 2.0 changes.

\begin{itemize}
\tightlist
\item
  \emph{Training.} All 35 task-seed fits completed; every one of the 175
  members gave finite predictions, so the retry rule was not used. 103
  members reached their lowest validation loss in the last three of the
  50 epochs (median epoch 48).
\item
  \emph{Where the results are.} M3 predictions are in
  \texttt{results/predictions\_m3/} and its analysis in
  \texttt{results/amendment2/}, apart from version 1\textquotesingle s
  files, which stay byte-identical
  (\texttt{results/revision/baseline\_manifest.json}). This separation
  was made before the M3 analysis was run.
\item
  \emph{Exploratory.} Written after M3 was scored and not part of the
  amendment: the ratio of mean spread on each held-out test to that on
  R, and the mean width of the C\_r-calibrated intervals, for each
  learner (\texttt{scripts/exploratory\_m3.py}).
\end{itemize}
